% Minimalne ustawienia pracy dyplomowej
% Zostawione głównie: marginesy, font, język, interlinia, nagłówki, podpisy, bibliografia i linki.

% Marginesy strony.
% inner = większy margines od strony oprawy przy druku dwustronnym.
\usepackage[
  top=2.5cm,
  bottom=2.5cm,
  inner=3.5cm,
  outer=2.5cm,
  a4paper,
  twoside
]{geometry}

% Obsługa grafik.
\usepackage{graphicx}

% Kodowanie fontów i pliku źródłowego — potrzebne przy kompilacji pdfLaTeX-em.
\usepackage[T1]{fontenc}
\usepackage[utf8]{inputenc}

% Język polski: dzielenie wyrazów, podpisy, nazwy elementów dokumentu.
\usepackage[polish]{babel}

% Font podobny do Times New Roman.
\usepackage{tgtermes}

% Formatowanie tytułów rozdziałów i sekcji.
\usepackage{titlesec}

% Formatowanie tytułu spisu treści.
\usepackage{tocloft}

% Interlinia, np. 1.5.
\usepackage{setspace}

% Wcięcie pierwszego akapitu po nagłówku.
\usepackage{indentfirst}

% Estetyczne tabele: \toprule, \midrule, \bottomrule.
\usepackage{booktabs}

% Formatowanie podpisów rysunków i tabel.
\usepackage[
  format=plain,
  labelfont=it,
  textfont=it,
  labelsep=period,
  justification=centering
]{caption}

% Podstawowe środowiska i symbole matematyczne.
\usepackage{amsmath}

% Nagłówki i stopki.
\usepackage{fancyhdr}

% Bibliografia.
\usepackage[style=ieee]{biblatex}
\usepackage{csquotes}

% Linki w PDF.
% hyperref najlepiej ładować pod koniec ustawień.
\usepackage{hyperref}
\hypersetup{
  colorlinks=true,
  citecolor=black,
  filecolor=black,
  linkcolor=black,
  urlcolor=black
}

% Ścieżka do grafik.
\graphicspath{{./images/}}

% Plik bibliografii.
\addbibresource{bibliografia.bib}

% Numerowanie rysunków i tabel w obrębie rozdziału, np. 2.1, 2.2.
\counterwithin{figure}{chapter}
\counterwithin{table}{chapter}

% Polska, krótka nazwa podpisu rysunku: "Rys." zamiast "Rysunek".
\addto\captionspolish{\renewcommand{\figurename}{Rys.}}

% Stopka z numerem strony na środku.
\pagestyle{fancy}
\fancyhf{}
\fancyfoot[C]{\fontsize{10}{12}\selectfont\thepage}
\renewcommand{\headrulewidth}{0pt}
\renewcommand{\footrulewidth}{0pt}

% Ten sam styl numeracji dla stron początkowych rozdziałów.
\fancypagestyle{plain}{
  \fancyhf{}
  \fancyfoot[C]{\fontsize{10}{12}\selectfont\thepage}
  \renewcommand{\headrulewidth}{0pt}
}

% Tytuł rozdziału: 12 pt, pogrubiony, wielkie litery.
\titleformat{\chapter}[block]
  {\normalfont\fontsize{12}{14}\selectfont\bfseries\MakeUppercase}
  {\thechapter.}{0.5em}{}
\titlespacing*{\chapter}{0pt}{12pt}{6pt}

% Tytuł sekcji: 10 pt, pogrubiony, kursywa.
\titleformat{\section}[block]
  {\normalfont\fontsize{10}{12}\selectfont\bfseries\itshape}
  {\thesection.}{0.5em}{}
\titlespacing*{\section}{0pt}{12pt}{6pt}

% Tytuł podsekcji: 10 pt, kursywa.
\titleformat{\subsection}[block]
  {\normalfont\fontsize{10}{12}\selectfont\itshape}
  {\thesubsection.}{0.5em}{}
\titlespacing*{\subsection}{0pt}{12pt}{6pt}

% Tytuł podpodsekcji: 10 pt, kursywa.
\titleformat{\subsubsection}[block]
  {\normalfont\fontsize{10}{12}\selectfont\itshape}
  {\thesubsubsection.}{0.5em}{}
\titlespacing*{\subsubsection}{0pt}{12pt}{6pt}

% Tytuł spisu treści: 12 pt, pogrubiony, wielkie litery.
\renewcommand{\cfttoctitlefont}{\fontsize{12}{14}\selectfont\bfseries\MakeUppercase}

% Wcięcie akapitu, brak dodatkowego odstępu między akapitami.
\setlength{\parindent}{1.25cm}
\setlength{\parskip}{0pt}

% Interlinia 1.5.
\onehalfspacing

% Pusta strona bez numeru — przydatne przy druku dwustronnym.
\newcommand\blankpage{%
  \null
  \thispagestyle{empty}%
  \addtocounter{page}{-1}%
  \newpage
}
